% ===== PRÉAMBULE CANONIQUE — à coller VERBATIM en tête de CHAQUE .tex =====
\documentclass[11pt,a4paper]{article}
\usepackage[utf8]{inputenc}\usepackage[T1]{fontenc}\usepackage{lmodern}
\usepackage[french]{babel}
\usepackage{amsmath,amssymb,mathtools}\usepackage{icomma}
\usepackage{siunitx}\sisetup{locale=FR}  % \qty{}{}, \si{}, \ang{}
\usepackage{xcolor}\usepackage{colortbl}
\usepackage{tikz}\usepackage{pgfplots}\pgfplotsset{compat=1.18}
\usepackage[siunitx,europeanresistors]{circuitikz} % circuits (élec.)
\usepackage[most]{tcolorbox}\usepackage{enumitem}\usepackage{array,booktabs}
\usepackage[a4paper,margin=2.2cm]{geometry}
\usetikzlibrary{arrows.meta,calc,angles,quotes,positioning,patterns,%
  backgrounds,shapes.geometric,decorations.pathmorphing,decorations.markings,intersections}
\definecolor{primary}{RGB}{222,49,99}\definecolor{primarydark}{RGB}{150,22,55}
\definecolor{defcol}{RGB}{0,114,178}\definecolor{methcol}{RGB}{52,92,148}
\definecolor{excol}{RGB}{0,135,90}\definecolor{attncol}{RGB}{200,40,55}
\tcbset{boxbase/.style={enhanced,breakable,boxrule=0.9pt,arc=2.5pt,
  left=3mm,right=3mm,top=2.3mm,bottom=2.3mm,fonttitle=\bfseries,coltitle=white}}
% --- boîtes TUTORIEL (noms EXACTS, ne pas renommer) ---
\newtcolorbox{cle}[1][]{boxbase,colback=defcol!6,colframe=defcol,
  title={\textsc{À retenir}\ifx\relax#1\relax\relax\else\textmd{ : \itshape #1}\fi}}
\newtcolorbox{methode}[1][]{boxbase,colback=methcol!7,colframe=methcol,sharp corners,
  title={\textsc{Méthode}\ifx\relax#1\relax\relax\else\textmd{ --- \itshape #1}\fi}}
\newtcolorbox{exemplebox}[1][]{boxbase,colback=excol!5,colframe=excol,
  title={\textsc{Exemple}\ifx\relax#1\relax\relax\else\textmd{ : \itshape #1}\fi}}
\newtcolorbox{piege}[1][]{boxbase,colback=attncol!6,colframe=attncol,
  title={\textsc{Piège}\ifx\relax#1\relax\relax\else\textmd{ : \itshape #1}\fi}}
\newtcolorbox{remarque}[1][]{boxbase,colback=black!4,colframe=black!45,
  title={\textsc{Remarque}\ifx\relax#1\relax\relax\else\textmd{ : \itshape #1}\fi}}
% --- boîtes EXERCICE / CORRIGÉ (noms EXACTS, auto-comptées) ---
\newtcolorbox[auto counter]{exo}[1][]{boxbase,colback=primary!4,colframe=primary,
  title={\textsc{Exercice}~\thetcbcounter\ifx\relax#1\relax\relax\else\textmd{ --- \itshape #1}\fi}}
\newtcolorbox[auto counter]{corrige}[1][]{boxbase,colback=excol!4,colframe=excol,
  title={\textsc{Corrigé}~\thetcbcounter\ifx\relax#1\relax\relax\else\textmd{ --- \itshape #1}\fi}}
\newcommand{\vv}[1]{\vec{#1}}\newcommand{\ds}{\displaystyle}
\newcommand{\R}{\mathbb{R}}\newcommand{\N}{\mathbb{N}}\newcommand{\Z}{\mathbb{Z}}
\begin{document}
% THEME_TITLE: Lentilles : familles, éléments et vergence
\sisetup{per-mode=symbol}

Une \textbf{lentille} est un bloc transparent limité par deux faces dont au moins une est courbe.
Tout se joue sur un seul critère visuel et une seule formule.

\begin{cle}[deux familles, un seul critère : l'épaisseur des bords]
\begin{itemize}[leftmargin=1.5em,topsep=2pt,itemsep=2pt]
  \item Lentille \textbf{convergente} : \textbf{bords minces} (épaisse au centre). Elle \emph{rassemble}
        les rayons. Distance focale \textcolor{excol}{$f>0$}. Symbole : segment à flèches vers
        l'\emph{extérieur}.
  \item Lentille \textbf{divergente} : \textbf{bords épais} (mince au centre). Elle \emph{écarte} les
        rayons. Distance focale \textcolor{attncol}{$f<0$}. Symbole : segment à flèches vers
        l'\emph{intérieur}.
\end{itemize}
Le \emph{nom} (biconvexe, ménisque\ldots) ne suffit pas : seul compte le rapport bord/centre.
\end{cle}

\begin{center}
\begin{tikzpicture}[scale=1.0]
  % convergente
  \begin{scope}
    \draw[black!35,dash pattern=on 3pt off 2pt] (-1.5,0)--(1.5,0);
    \fill[defcol!25] plot[smooth cycle] coordinates {(0,1) (0.24,0) (0,-1) (-0.24,0)};
    \draw[defcol!80,line width=0.5pt] plot[smooth cycle] coordinates {(0,1) (0.24,0) (0,-1) (-0.24,0)};
    \node[font=\scriptsize,align=center] at (0,-1.45) {bords \textbf{minces}\\ $\Rightarrow$ \textcolor{excol}{convergente}};
  \end{scope}
  % symbole convergente
  \begin{scope}[shift={(3.0,0)}]
    \draw[black!35,dash pattern=on 3pt off 2pt] (-1.2,0)--(1.2,0);
    \draw[line width=1.6pt,<->,>=Stealth,black!65] (0,-1)--(0,1);
    \node[font=\scriptsize] at (0,-1.45) {symbole};
  \end{scope}
  % divergente
  \begin{scope}[shift={(6.5,0)}]
    \draw[black!35,dash pattern=on 3pt off 2pt] (-1.5,0)--(1.5,0);
    \fill[defcol!25] (-0.32,-1)--(0.32,-1)--(0.08,0)--(0.32,1)--(-0.32,1)--(-0.08,0)--cycle;
    \draw[defcol!80,line width=0.5pt] (-0.32,-1)--(0.32,-1)--(0.08,0)--(0.32,1)--(-0.32,1)--(-0.08,0)--cycle;
    \node[font=\scriptsize,align=center] at (0,-1.45) {bords \textbf{épais}\\ $\Rightarrow$ \textcolor{attncol}{divergente}};
  \end{scope}
  % symbole divergente
  \begin{scope}[shift={(9.5,0)}]
    \draw[black!35,dash pattern=on 3pt off 2pt] (-1.2,0)--(1.2,0);
    \draw[line width=1.6pt,>-<,>=Stealth,black!65] (0,-1)--(0,1);
    \node[font=\scriptsize] at (0,-1.45) {symbole};
  \end{scope}
\end{tikzpicture}
\end{center}

\begin{cle}[les éléments d'une lentille mince]
\begin{itemize}[leftmargin=1.5em,topsep=2pt,itemsep=1pt]
  \item Le \textbf{centre optique} $O$ : tout rayon qui le traverse ressort \textbf{sans être dévié}.
  \item Le \textbf{foyer image} $F'$ : les rayons arrivant \emph{parallèles à l'axe} convergent en $F'$.
  \item Le \textbf{foyer objet} $F$ : symétrique de $F'$ par rapport à $O$.
  \item La \textbf{distance focale} $f=\overline{OF'}$ : $F$ et $F'$ sont à la distance $|f|$ de $O$.
\end{itemize}
\end{cle}

\begin{center}
\begin{tikzpicture}[scale=0.88,>=Stealth]
  \draw[black!35,dash pattern=on 3pt off 2pt] (-4.4,0)--(4.4,0);
  \draw[line width=1.6pt,<->,black!65] (0,-1.5)--(0,1.5);
  \fill[black] (0,0) circle (1.4pt); \node[below right,font=\footnotesize] at (0,0) {$O$};
  \fill[primary] (2.6,0) circle (1.6pt); \node[below,font=\footnotesize] at (2.6,0) {$F'$};
  \fill[primary] (-2.6,0) circle (1.6pt); \node[below,font=\footnotesize] at (-2.6,0) {$F$};
  % rayon par O non devie
  \draw[black!75,line width=0.9pt,->] (-4.0,1.0)--(0,0);
  \draw[defcol,line width=0.9pt,->] (0,0)--(3.8,-0.95);
  \node[black!70,above,font=\scriptsize] at (-2.0,0.6) {par $O$ : non dévié};
  \draw[<->,primary,line width=0.6pt] (0,1.15)--(2.6,1.15);
  \node[primary,above,font=\scriptsize] at (1.3,1.15) {$f=\overline{OF'}$};
\end{tikzpicture}
\end{center}

\begin{cle}[vergence $C=\dfrac{1}{f}$ --- la formule reine]
La \textbf{vergence} mesure la puissance d'une lentille :
\[ \boxed{\,C=\dfrac{1}{f}\,}\qquad f \text{ \textbf{en mètres}},\qquad
   \begin{cases}C>0 & \text{convergente}\\ C<0 & \text{divergente}\end{cases} \]
Elle s'exprime en \textbf{dioptries} ($\delta$), avec $1~\delta = 1~\mathrm{m^{-1}}$. Plus $|C|$ est
\emph{grand}, plus la lentille est \emph{puissante} (petit $|f|$).
\end{cle}

\begin{cle}[lentilles accolées : les vergences s'ajoutent]
Plusieurs lentilles minces \textbf{accolées} (centres confondus) équivalent à une seule lentille de
vergence
\[ \boxed{\,C_{\text{éq}}=C_1+C_2+\cdots=\sum_k C_k\,}\qquad
   \Big(\text{soit } \tfrac{1}{f_{\text{éq}}}=\textstyle\sum_k \tfrac{1}{f_k}\Big). \]
\end{cle}

\begin{piege}[trois erreurs à ne jamais commettre]
\begin{itemize}[leftmargin=1.5em,topsep=2pt,itemsep=1pt]
  \item La vergence est l'\textbf{inverse} de $f$, \emph{pas} son opposé : $C=1/f$, jamais $C=-f$.
  \item $f$ doit être \textbf{en mètres} avant tout calcul. Une lentille de $f=\qty{20}{\cm}$ a
        $C=\tfrac{1}{0,20}=+5~\delta$, et surtout \emph{pas} $\tfrac{1}{20}$.
  \item Une \textbf{loupe est convergente} (bords minces, $f>0$), jamais divergente.
\end{itemize}
\end{piege}

\begin{exemplebox}[de $f$ à $C$ et retour]
Lentille $f=+\qty{25}{\cm}=+\qty{0.25}{\m}$ : $C=\dfrac{1}{0,25}=+4~\delta$ (convergente).\\
Verre correcteur $C=-2~\delta$ : $f=\dfrac{1}{-2}=-\qty{0.5}{\m}=-\qty{50}{\cm}$ (divergente).\\
On accole ces deux lentilles : $C_{\text{éq}}=+4-2=+2~\delta$, encore \emph{convergente},
$f_{\text{éq}}=\dfrac{1}{2}=+\qty{0.5}{\m}$.
\end{exemplebox}

\end{document}
